\section{Introduction}
\label{sec:intro}

Neural Radiance Fields (NeRFs) have consolidated themselves in recent years as the state-of-the-art for novel view synthesis, in both industry and academia \cite{kerbl3Dgaussians, Tancik2023, substance3dviewer}. However, the most advanced neural rendering techniques require considerable computational resources and multi-view datasets for training. Both challenges make practical applications, such as mobile scanning, virtual try-ons, and synthesis in reduced datasets, unfeasible.To address this problem, several methods have been proposed to advance the state of the art, including Instant-NGP (\textit{Instant Neural Graphics Primitives}), Nerfacto, PixelNeRF, and GNT (\textit{Generalizable NeRF Transformer}) \cite{muller2022, Tancik2023, yu2021pixelnerfneuralradiancefields, t2023attentionnerfneeds}. These methods investigate different aspects of neural rendering, such as spatial \textit{encoding} \cite{muller2022}, hybrid architectures that combine Convolutional Neural Networks (CNNs) and Vision Transformers (ViTs) \cite{schwarz2021grafgenerativeradiancefields, wu2021cvt, t2023attentionnerfneeds}, and model pre-computation (\textit{baking}) into compact representations with a reduced memory footprint \cite{hedman2021bakingneuralradiancefields}.

In this context, this work investigates the application of different NeRF models and their respective optimization techniques to sparse-view datasets by developing a methodology for evaluating the trade-off between computational performance and the visual quality of the reconstructed volumes. Reconstruction fidelity was assessed using three widely adopted image quality metrics: PSNR (\textit{Peak Signal-to-Noise Ratio}), SSIM (\textit{Structural Similarity Index Measure}) \cite{wang2004image}, and LPIPS (\textit{Learned Perceptual Image Patch Similarity}) \cite{zhang2018unreasonable}. The operational cost of each model was quantified using system-level performance metrics, including training convergence time, inference speed, and memory footprint, particularly VRAM (\textit{Video Random Access Memory}) consumption. Although previous studies have investigated NeRFs under limited-data settings \cite{niemeyer2022regnerf, kim2023poserefinement}, and benchmarking frameworks such as NerfBaselines \cite{kulhanek2025nerfbaselinesconsistentreproducibleevaluation} have evaluated traditional and per-scene NeRF models, a systematic comparison of generalizable neural rendering architectures employing convolutional and attention-based feature extractors under controlled hardware conditions and reproducible data splits remains absent from the literature. To bridge this gap, this work develops a systematic evaluation methodology for quantitatively assessing the trade-offs between computational efficiency and reconstruction fidelity across four neural rendering architectures under three sparse-view data regimes. Thus, the main contributions of this work are: (i) the formulation of a unified evaluation protocol for few-shot and per-scene NeRF architectures, ensuring reproducibility through standardized data partitions, controlled execution environments, and consistent evaluation metrics across three sparse-view data regimes; and (ii) the implementation and public release of integration wrappers that extend the Nerfstudio framework with PixelNeRF and GNT, providing a common experimental infrastructure for standardized benchmarking of heterogeneous neural rendering architectures.

